\renewcommand{\thetheorem}{16.\arabic{theorem}}
\setcounter{theorem}{0}

\setcounter{chapter}{15} % Sets chapter counter so the next chapter is 16

% ---------------------------------------------------------------------------
% Provenance of the prose in this file.
%   % Extractor: <model>   the block below was transcribed from Prof. Biran's
%                          handwritten notes by that model
%   % Creator:   <model>   the block below was written by that model and has no
%                          counterpart in the notes
% A tag holds until the next tag.
% ---------------------------------------------------------------------------

\chapter{Linear Maps and Bases}

% Extractor: Gemini 3.1 Pro
\subsection{Writing Linear Maps Using Bases and Coordinates}


In this lecture, we will assume all vector spaces to be \qt{finite-di\-men\-sional} unless otherwise stated. We will write bases of a vector space $V$ as an ordered list (or tuple) $\mathcal{B} = (v_1, \dots, v_n)$.

% def:16.1
\begin{definition}[Coordinate Vector]
\label{def:coordinate_vector}
Let $V$ be a vector space over $K$, and let $\mathcal{B} = (v_1, \dots, v_n)$ be a basis of $V$. Let $v \in V$. We define the \newterm{coordinate vector} of $v$ with respect to the basis $\mathcal{B}$ as
\[
[v]_{\mathcal{B}} := 
\begin{pmatrix}
a_1 \\
\vdots \\
a_n
\end{pmatrix}
\in K^n,
\]
where $a_1, \dots, a_n \in K$ are the unique scalars in $K$ such that $v = a_1 v_1 + \dots + a_n v_n$.
\end{definition}

\begin{remark}
Note that $[v]_{\mathcal{B}}$ is uniquely defined, since every $v \in V$ has a unique representation as a linear combination of the basis vectors $v_1, \dots, v_n$.
\end{remark}

Define a map $\Phi_{\mathcal{B}}: V \to K^n$ given by:
\[
\Phi_{\mathcal{B}}(v) := [v]_{\mathcal{B}} \in K^n.
\]

% prop:16.2
\begin{proposition}[Isomorphism of the Coordinate Map]
\label{prop:coordinate_map_isomorphism}
The map $\Phi_{\mathcal{B}}$ is an isomorphism.
\end{proposition}

\begin{proof}
We first show that $\Phi_{\mathcal{B}}$ is a linear map. Let $a, b \in K$, and let $v, v' \in V$. Write $v$ and $v'$ as linear combinations of the elements of $\mathcal{B}$:
\[
v = a_1 v_1 + \dots + a_n v_n, \quad \text{and} \quad v' = b_1 v_1 + \dots + b_n v_n,
\]
with $a_i, b_i \in K$. By definition, we have:
\[
[v]_{\mathcal{B}} = 
\begin{pmatrix}
a_1 \\
\vdots \\
a_n
\end{pmatrix},
\quad \text{and} \quad
[v']_{\mathcal{B}} = 
\begin{pmatrix}
b_1 \\
\vdots \\
b_n
\end{pmatrix}.
\]
Now, $a v + b v' = (a a_1 + b b_1) v_1 + \dots + (a a_n + b b_n) v_n$, and hence:
\[
\Phi_{\mathcal{B}}(a v + b v') = [a v + b v']_{\mathcal{B}} = 
\begin{pmatrix}
a a_1 + b b_1 \\
\vdots \\
a a_n + b b_n
\end{pmatrix}
= a 
\begin{pmatrix}
a_1 \\
\vdots \\
a_n
\end{pmatrix}
+ b
\begin{pmatrix}
b_1 \\
\vdots \\
b_n
\end{pmatrix}
= a \Phi_{\mathcal{B}}(v) + b \Phi_{\mathcal{B}}(v').
\]
This proves that $\Phi_{\mathcal{B}}$ is linear. 

We now claim that $\Phi_{\mathcal{B}}$ is an isomorphism. By previous results (\cref{lem:bijective_is_iso}), it is enough to show $\Phi_{\mathcal{B}}$ is bijective. Indeed, if $v \in \ker(\Phi_{\mathcal{B}})$, then $v = 0 v_1 + \dots + 0 v_n = 0$, which implies $\ker(\Phi_{\mathcal{B}}) = \{0\}$. This shows that $\Phi_{\mathcal{B}}$ is injective. 

For the surjectivity of $\Phi_{\mathcal{B}}$, let $(a_1, \dots, a_n)\transp \in K^n$. Put $v := a_1 v_1 + \dots + a_n v_n \in V$. Then, $\Phi_{\mathcal{B}}(v) = (a_1, \dots, a_n)\transp$, which implies $\Img(\Phi_{\mathcal{B}}) = K^n$.
\end{proof}

\begin{remark}
The inverse map to $\Phi_{\mathcal{B}}$ is given by:
\[
\Phi_{\mathcal{B}}^{-1}: K^n \to V, \quad \Phi_{\mathcal{B}}^{-1}
\begin{pmatrix}
a_1 \\
\vdots \\
a_n
\end{pmatrix}
:= a_1 v_1 + \dots + a_n v_n.
\]
\end{remark}

\begin{remark}
The previous \cref{prop:coordinate_map_isomorphism} gives yet another proof for the following Corollary.
\end{remark}

% cor:16.3
\begin{corollary}[Classification of Finite-Dimensional Vector Spaces]
\label{cor:vs_isomorphic_to_kn}
Let $V$ be a vector space over $K$ with $\dim V = n \in \mathbb{Z}_{\ge 0}$. Then, $V \cong K^n$.
\end{corollary}

\begin{proof}
Choose a basis $\mathcal{B}$ for $V$. Then, $\Phi_{\mathcal{B}}: V \to K^n$ is an isomorphism.
\end{proof}

\begin{remark}
\begin{itemize}
    \item The Corollary does NOT hold for $\dim V = \infty$, at least not as stated.
    \item We have seen that $V \cong K^n$ where $n = \dim V < \infty$, but in general, there does not exist a \qt{canonical} (or preferred) isomorphism $V \to K^n$. For example, $\Phi_{\mathcal{B}}$ depends on a choice of a basis $\mathcal{B}$.
\end{itemize}
\end{remark}

% ex:16.a
\begin{example}[Coordinate Vector Calculation]
\label{ex:coordinate_vector_calculation}
Consider $V = \mathbb{R}^2$, let $\mathcal{E} = \left(\binom{1}{0}, \binom{0}{1}\right)$ be the standard basis, and let $\mathcal{B} = \left(\binom{1}{1}, \binom{1}{-1}\right)$ be another basis (check that $\mathcal{B}$ is indeed a basis!). We have
\[
\left[\binom{x}{y}\right]_{\mathcal{E}} = \binom{x}{y}, \quad \text{and we want to find} \quad \left[\binom{x}{y}\right]_{\mathcal{B}} = \binom{a}{b} = ?
\]
Let us first try to determine $\left[\binom{1}{0}\right]_{\mathcal{B}}$ and $\left[\binom{0}{1}\right]_{\mathcal{B}}$. We set $\left[\binom{1}{0}\right]_{\mathcal{B}} = \binom{a_1}{a_2}$, where
\[
a_1 \begin{pmatrix} 1 \\ 1 \end{pmatrix} + a_2 \begin{pmatrix} 1 \\ -1 \end{pmatrix} = \begin{pmatrix} 1 \\ 0 \end{pmatrix}.
\]
This leads to the system of linear equations:
\[
\begin{cases}
a_1 + a_2 = 1 \\
a_1 - a_2 = 0
\end{cases}
\]
Solving this system yields $a_1 = a_2 = \frac{1}{2}$, and therefore, $\left[\binom{1}{0}\right]_{\mathcal{B}} = \binom{1/2}{1/2}$.

Similarly, for the second standard basis vector, we set
\[
\left[\binom{0}{1}\right]_{\mathcal{B}} = \begin{pmatrix} b_1 \\ b_2 \end{pmatrix}, \quad \text{where} \quad b_1 \begin{pmatrix} 1 \\ 1 \end{pmatrix} + b_2 \begin{pmatrix} 1 \\ -1 \end{pmatrix} = \begin{pmatrix} 0 \\ 1 \end{pmatrix}.
\]
The resulting system is:
\[
\begin{cases}
b_1 + b_2 = 0 \\
b_1 - b_2 = 1
\end{cases}
\]
Solving this yields $b_1 = \frac{1}{2}$ and $b_2 = -\frac{1}{2}$, so $\left[\binom{0}{1}\right]_{\mathcal{B}} = \binom{1/2}{-1/2}$.

Finally, using the fact that $\Phi_{\mathcal{B}}$ is linear, we calculate:
\[
\left[\binom{x}{y}\right]_{\mathcal{B}} = x \left[\binom{1}{0}\right]_{\mathcal{B}} + y \left[\binom{0}{1}\right]_{\mathcal{B}} = \begin{pmatrix} x/2 \\ x/2 \end{pmatrix} + \begin{pmatrix} y/2 \\ -y/2 \end{pmatrix} = \begin{pmatrix} \frac{x+y}{2} \\ \frac{x-y}{2} \end{pmatrix}.
\]
\end{example}

\subsection{Representation Matrices}

Now we arrive at the main question. Let $V$ and $W$ be vector spaces over $K$, let $n := \dim V$ and $m := \dim W$, and let $T: V \to W$ be a linear map. Fix a basis $\mathcal{B}$ for $V$ and a basis $\mathcal{C}$ for $W$. Consider the following diagram, in which the two paths from $V$ to $K^m$ agree, so that the diagram is commutative:
\begin{center}
\begin{tikzcd}[row sep=large, column sep=large]
V \arrow[r, "T"] \arrow[d, "\Phi_{\mathcal{B}}"'] & W \arrow[d, "\Phi_{\mathcal{C}}"] \\
K^n \arrow[r, "\Phi_{\mathcal{C}} \circ T \circ \Phi_{\mathcal{B}}^{-1}"'] & K^m
\end{tikzcd}
\end{center}
Since $\Phi_{\mathcal{B}}^{-1}$, $T$, and $\Phi_{\mathcal{C}}$ are all linear maps, the map $\Phi_{\mathcal{C}} \circ T \circ \Phi_{\mathcal{B}}^{-1}: K^n \to K^m$ is also linear. By previous results (\cref{lem:all_linear_maps_are_matrices}), any linear map $K^n \to K^m$ can be written using a matrix $A \in \M_{m \times n}(K)$. In other words, there exists a unique matrix $A \in \M_{m \times n}(K)$ such that $\Phi_{\mathcal{C}} \circ T \circ \Phi_{\mathcal{B}}^{-1} = T_A$.

The matrix $A$ is called the \newterm{representation matrix} of $T$ with respect to the bases $\mathcal{B}$ and $\mathcal{C}$. We also say that $A$ represents $T$ in the bases $\mathcal{B}$ and $\mathcal{C}$. We denote this matrix by $[T]_{\mathcal{C}}^{\mathcal{B}} \in \M_{m \times n}(K)$.

\begin{question}
\begin{enumerate}
    \item How can we calculate the entries of $A$?
    \item How does $A$ depend on different choices of the bases $\mathcal{B}$ and $\mathcal{C}$?
    \item If we have three vector spaces $V, W, U$ with bases $\mathcal{B}, \mathcal{C}, \mathcal{D}$ respectively, and $T: V \to W$ and $S: W \to U$ are linear maps where $T$ is represented by the matrix $M$ and $S$ by the matrix $N$, then what is the matrix representing $S \circ T: V \to U$ with respect to the bases $\mathcal{B}$ and $\mathcal{D}$?
\end{enumerate}
\end{question}

Let $T: V \to W$ be a linear map, let $\mathcal{B} = (v_1, \dots, v_n)$ be a basis for $V$, and let $\mathcal{C} = (w_1, \dots, w_m)$ be a basis for $W$. We wish to determine the entries of its representation matrix $A := [T]_{\mathcal{C}}^{\mathcal{B}} \in \M_{m \times n}(K)$. By definition, $A$ is the unique matrix such that the linear map $T_A: K^n \to K^m$, given by left-multiplication by $A$, satisfies $T_A = \Phi_{\mathcal{C}} \circ T \circ \Phi_{\mathcal{B}}^{-1}$. 

Recall that the $j$-th column of $A$ is given by $T_A(e_j)$, where $e_j \in K^n$ is the $j$-th standard basis vector. We can evaluate $T_A(e_j)$ using the commutative diagram:
\[
T_A(e_j) = \left( \Phi_{\mathcal{C}} \circ T \circ \Phi_{\mathcal{B}}^{-1} \right)(e_j) = \Phi_{\mathcal{C}}\left( T\left( \Phi_{\mathcal{B}}^{-1}(e_j) \right) \right).
\]
Recall that the coordinate vector of the basis vector $v_j \in V$ with respect to its own basis $\mathcal{B}$ is the $j$-th standard basis vector $e_j \in K^n$, meaning $\Phi_{\mathcal{B}}(v_j) = [v_j]_{\mathcal{B}} = e_j$. Applying the inverse coordinate map gives $\Phi_{\mathcal{B}}^{-1}(e_j) = v_j$. Substituting this into our expression, we have:
\[
T_A(e_j) = \Phi_{\mathcal{C}}(T(v_j)) = [T(v_j)]_{\mathcal{C}}.
\]
This means that the $j$-th column of the representation matrix $[T]_{\mathcal{C}}^{\mathcal{B}}$ is precisely the coordinate vector of $T(v_j)$ with respect to the basis $\mathcal{C}$. Thus, we can write $A$ in block-column form as:
\[
A = [T]_{\mathcal{C}}^{\mathcal{B}} = 
\begin{pmatrix}
| & & | \\
[T(v_1)]_{\mathcal{C}} & \dots & [T(v_n)]_{\mathcal{C}} \\
| & & |
\end{pmatrix}
\in \M_{m \times n}(K).
\]
Equivalently, if we write $A = (a_{ij}) \in \M_{m \times n}(K)$, the entries $a_{ij}$ of the $j$-th column are determined by expressing $T(v_j) \in W$ as a linear combination of the basis vectors $\mathcal{C} = (w_1, \dots, w_m)$:
\[
T(v_j) = a_{1j} w_1 + a_{2j} w_2 + \dots + a_{mj} w_m = \sum_{i=1}^m a_{ij} w_i.
\]

% prop:16.4
\begin{proposition}[Evaluation via Representation Matrix]
\label{prop:matrix_evaluation}
Let $T: V \to W$ be a linear map between two finite-di\-men\-sional vector spaces $V$ and $W$. Let $\mathcal{B}$ and $\mathcal{C}$ be bases for $V$ and $W$, respectively. Then, for all $v \in V$:
\[
[T]_{\mathcal{C}}^{\mathcal{B}} \cdot [v]_{\mathcal{B}} = [T(v)]_{\mathcal{C}}.
\]
\end{proposition}

\begin{proof}
% Creator: Gemini 3.1 Pro , the notes give only the structural proof below
\textbf{First Proof (Computational):} Let $v \in V$ be an arbitrary vector. We can express $v$ uniquely as a linear combination of the basis vectors in $\mathcal{B} = (v_1, \dots, v_n)$:
\[
v = \sum_{j=1}^n x_j v_j.
\]
By definition, the coordinate vector of $v$ with respect to $\mathcal{B}$ is $[v]_{\mathcal{B}} = (x_1, \dots, x_n)\transp \in K^n$. Since $T$ is linear, we can apply it to $v$:
\[
T(v) = T\left( \sum_{j=1}^n x_j v_j \right) = \sum_{j=1}^n x_j T(v_j).
\]
Now, we take the coordinate vector of $T(v)$ with respect to the basis $\mathcal{C}$. Since the coordinate map $\Phi_{\mathcal{C}}$ is an isomorphism, it is linear, and we obtain:
\[
[T(v)]_{\mathcal{C}} = \Phi_{\mathcal{C}}(T(v)) = \Phi_{\mathcal{C}}\left( \sum_{j=1}^n x_j T(v_j) \right) = \sum_{j=1}^n x_j [T(v_j)]_{\mathcal{C}}.
\]
Observe that this last expression is exactly the definition of the matrix-vector product of the matrix $[T]_{\mathcal{C}}^{\mathcal{B}}$ (whose columns are $[T(v_j)]_{\mathcal{C}}$) and the column vector $[v]_{\mathcal{B}}$ (whose entries are $x_j$):
\[
\sum_{j=1}^n x_j [T(v_j)]_{\mathcal{C}} = 
\begin{pmatrix}
| & & | \\
[T(v_1)]_{\mathcal{C}} & \dots & [T(v_n)]_{\mathcal{C}} \\
| & & |
\end{pmatrix}
\begin{pmatrix}
x_1 \\ \vdots \\ x_n
\end{pmatrix}
= [T]_{\mathcal{C}}^{\mathcal{B}} \cdot [v]_{\mathcal{B}}.
\]
Therefore, we conclude that $[T]_{\mathcal{C}}^{\mathcal{B}} \cdot [v]_{\mathcal{B}} = [T(v)]_{\mathcal{C}}$.

% Extractor: Gemini 3.1 Pro
\textbf{Second Proof (Structural):} Recall that the representation matrix $A := [T]_{\mathcal{C}}^{\mathcal{B}}$ is uniquely defined by the relation $T_A = \Phi_{\mathcal{C}} \circ T \circ \Phi_{\mathcal{B}}^{-1}$. Evaluating the linear transformation $T_A$ at the coordinate vector $[v]_{\mathcal{B}} = \Phi_{\mathcal{B}}(v) \in K^n$ yields:
\[
T_A([v]_{\mathcal{B}}) = \left(\Phi_{\mathcal{C}} \circ T \circ \Phi_{\mathcal{B}}^{-1}\right) (\Phi_{\mathcal{B}}(v)) = \Phi_{\mathcal{C}}\left(T\left(\Phi_{\mathcal{B}}^{-1}(\Phi_{\mathcal{B}}(v))\right)\right) = \Phi_{\mathcal{C}}(T(v)) = [T(v)]_{\mathcal{C}}.
\]
On the other hand, the map $T_A$ acts on vectors in $K^n$ by left-multiplication by the matrix $A$. Thus, we also have:
\[
T_A([v]_{\mathcal{B}}) = A \cdot [v]_{\mathcal{B}} = [T]_{\mathcal{C}}^{\mathcal{B}} \cdot [v]_{\mathcal{B}}.
\]
Equating both expressions gives $[T]_{\mathcal{C}}^{\mathcal{B}} \cdot [v]_{\mathcal{B}} = [T(v)]_{\mathcal{C}}$, which completes the proof.
\end{proof}

\begin{example}[Examples of Representation Matrices]
\begin{enumerate}
    \item  Let $V = K^n$ and $W = K^m$, and let $T: K^n \to K^m$ be given by $T = T_A$ for $A \in \M_{m \times n}(K)$, so $T(v) = A \cdot v$. Take $\mathcal{B} := (e_1, \dots, e_n)$ and $\mathcal{C} := (e_1, \dots, e_m)$ to be the standard bases. Then, $[T]_{\mathcal{C}}^{\mathcal{B}} = A$. \textit{(Check the details!)}
    \item  Let $V = K[x]_3$ and $W = K[x]_2$, and let $D: V \to W$ be the derivative map defined by $D(p(x)) = p'(x)$. Let $\mathcal{B} := (1, x, x^2, x^3)$ and $\mathcal{C} := (1, x, x^2)$. We calculate:
    \begin{align*}
    D(1) &= 0 = 0 \cdot 1 + 0 \cdot x + 0 \cdot x^2 \\
    D(x) &= 1 = 1 \cdot 1 + 0 \cdot x + 0 \cdot x^2 \\
    D(x^2) &= 2x = 0 \cdot 1 + 2 \cdot x + 0 \cdot x^2 \\
    D(x^3) &= 3x^2 = 0 \cdot 1 + 0 \cdot x + 3 \cdot x^2
    \end{align*}
    Therefore, the representation matrix is:
    \[
    [D]_{\mathcal{C}}^{\mathcal{B}} = 
    \begin{pmatrix}
    0 & 1 & 0 & 0 \\
    0 & 0 & 2 & 0 \\
    0 & 0 & 0 & 3
    \end{pmatrix}.
    \]
    \item  Revisiting the previous item, suppose we view $D$ as a linear map $D: K[x]_3 \to K[x]_3$, so that the same map is now given the same basis on both sides. Then,
    \[
    [D]_{\mathcal{B}}^{\mathcal{B}} = 
    \begin{pmatrix}
    0 & 1 & 0 & 0 \\
    0 & 0 & 2 & 0 \\
    0 & 0 & 0 & 3 \\
    0 & 0 & 0 & 0
    \end{pmatrix}.
    \]
\end{enumerate}
\end{example}


\subsection{Representing Composition of Linear Maps by Matrices}

Let $T: V \to W$ and $S: W \to U$ be linear maps. Let $\mathcal{A} = (v_1, \dots, v_n)$ be a basis for $V$, let $\mathcal{B} = (w_1, \dots, w_m)$ be a basis for $W$, and let $\mathcal{C} = (u_1, \dots, u_p)$ be a basis for $U$. The dimensions are thus $n = \dim V$, $m = \dim W$, and $p = \dim U$.

Consider the composition $S \circ T: V \longrightarrow U$.

\begin{question}
What is the matrix $[S \circ T]_{\mathcal{C}}^{\mathcal{A}}$? How is it related to the matrices $[T]_{\mathcal{B}}^{\mathcal{A}}$ and $[S]_{\mathcal{C}}^{\mathcal{B}}$?
\end{question}

To answer this, let us write the representation matrices as follows:
\[
[T]_{\mathcal{B}}^{\mathcal{A}} = (a_{ij}), \quad [S]_{\mathcal{C}}^{\mathcal{B}} = (b_{ij}), \quad [S \circ T]_{\mathcal{C}}^{\mathcal{A}} = (c_{ij}).
\]
By definition, the images of the basis vectors are given by:
\begin{equation}
\label{eq:T_basis_image}
T(v_j) = \sum_{k=1}^m a_{kj} w_k, \quad \forall \, 1 \le j \le n,
\end{equation}
\begin{equation}
\label{eq:S_basis_image}
S(w_k) = \sum_{i=1}^p b_{ik} u_i, \quad \forall \, 1 \le k \le m,
\end{equation}
\begin{equation}
\label{eq:composition_basis_image}
(S \circ T)(v_j) = \sum_{i=1}^p c_{ij} u_i, \quad \forall \, 1 \le j \le n.
\end{equation}

But we can also calculate $(S \circ T)(v_j)$ using \eqref{eq:T_basis_image} and \eqref{eq:S_basis_image}:
\begin{align*}
(S \circ T)(v_j) &= S \left( \sum_{k=1}^m a_{kj} w_k \right) \\
&= \sum_{k=1}^m a_{kj} S(w_k) \\
&= \sum_{k=1}^m a_{kj} \left( \sum_{i=1}^p b_{ik} u_i \right) \\
&= \sum_{k=1}^m \sum_{i=1}^p a_{kj} b_{ik} u_i = \sum_{i=1}^p \left( \sum_{k=1}^m b_{ik} a_{kj} \right) \cdot u_i.
\end{align*}
Comparing this with \eqref{eq:composition_basis_image}, we find that $c_{ij} = \sum_{k=1}^m b_{ik} a_{kj}$. In other words, row $i$ of $[S]_{\mathcal{C}}^{\mathcal{B}}$ multiplied by column $j$ of $[T]_{\mathcal{B}}^{\mathcal{A}}$ yields the entry $c_{ij}$:
\[
\begin{pmatrix}
b_{i1} & \dots & b_{im}
\end{pmatrix}
\cdot 
\begin{pmatrix}
a_{1j} \\
\vdots \\
a_{mj}
\end{pmatrix}
= c_{ij}.
\]
This reveals that:
\[
[S]_{\mathcal{C}}^{\mathcal{B}} \cdot [T]_{\mathcal{B}}^{\mathcal{A}} = [S \circ T]_{\mathcal{C}}^{\mathcal{A}},
\]
where the matrix dimensions follow $(p \times m) \cdot (m \times n) = (p \times n)$.

In order to obtain $[S \circ T]_{\mathcal{C}}^{\mathcal{A}}$ from $[T]_{\mathcal{B}}^{\mathcal{A}}$ and $[S]_{\mathcal{C}}^{\mathcal{B}}$, we need to use a new operation on matrices. It is called \qt{matrix multiplication}. But is this operation really new?

\medskip
\noindent--- Yes, but not entirely\dots

% Creator: Opus 5 (High)
\begin{ainote}
This chapter came through the transcription essentially intact: every statement
of the notes is present, in the notes' order, and the printed numbers
16.1--16.4 already agreed with the handwritten ones. The abrupt final line is
Prof. Biran's own; his notes end there, in the middle of the thought that
chapter 17 picks up.

What needed repair was small and mostly typographic. The third example carried a
\verb|\setcounter{enumi}{1}|, annotated as matching the professor's numbering,
which made the list print two items labelled \textbf{(b)}; the notes label that
item $2'$ because it revisits the previous one, so the counter manipulation is
dropped and the revisiting is said in words instead. The sentence introducing
the commutative diagram ended in the word \qt{test}, evidently a leftover. A
\verb|\usetikzlibrary{babel}| sat in the middle of the document, although the
preamble already loads that library. Prof. Biran's aside \qt{check the
details!} on the first representation-matrix example is restored, and his
closing line is now set as the reply it is rather than as a stray hyphen.

The proof that $\Phi_{\mathcal{B}}$ is an isomorphism cited the definition of
isomorphic spaces for the step that bijectivity suffices. That step is
\cref{lem:bijective_is_iso}, which the transcript of the previous chapter had
dropped altogether and which is now restored; the citation points there.

One addition of the transcript is worth keeping and worth flagging: the notes
prove \cref{prop:matrix_evaluation} only structurally, by observing that the
asserted identity is exactly the commutative diagram defining
$[T]_{\mathcal{C}}^{\mathcal{B}}$. The computational proof that precedes it in
the text above is not Prof. Biran's.
\end{ainote}
